\paragraph{A. Proof of the fixed-set Lemma 19.}
Begin with $R$ equal to the uniform distribution over $W^s$. If there are $q$ words and $q$ values such that the probability of getting those values in $B$ is larger than $(1+\alpha)/W^q$ then we fix them to those values, in both $B$ and $R$. Now we have subsets of $W^{s-q}$, the loss has decreased by an additive $\log_2 1/(1+\alpha)=\Omega(\alpha)$, and we repeat the process.
Because the initial loss was $b$, this process stops after $O(b/\alpha)$ iterations. In the end, the loss inside the final universe is at most $b$, since we never increase loss. With respect to the original universe, because we fixed $O(qb/\alpha)$ words, the loss is at most $O(wqb/\alpha)+b\leq b\cdot O(wq/\alpha)$.
Let $B_1$ and $R$ be the distributions when the process stops. Consider any tree $g:W^s\to\{0,1\}$ of depth $q$. Let $P$ be the collection of paths in $g$ leading to the output 1. Note that each path $p\in P$ corresponds to $q$ input words and $q$ values for them. Write $P_X[p]$ for the probability of following path $p$ under distribution $X$. By above we have $P_{B_1}[p]\leq(1+\alpha)/W^q=(1+\alpha)P_R[p]$.
Hence
\[\mathbb P[g(B_1)=1]=\sum_{p\in P}P_{B_1}[p]\leq\sum_{p\in P}(1+\alpha)P_R[p]=(1+\alpha)\mathbb P[g(R)=1].\]
And so in particular $\mathbb P[g(B_1)=1]\leq\mathbb P[g(R)=1]+\alpha$.
Repeating the argument with 0 and 1 swapped yields the lemma for trees with boolean alphabet. To prove the lemma for a tree $g'$ with arbitrary alphabet, reduce to the case of boolean alphabet in the following standard way. Suppose that the statistical distance between $g'(R)$ and $g'(B_1)$ is $>\alpha$. This means that there exists a set $T$ such that
\[|\mathbb P[g'(R)\in T]-\mathbb P[g'(B_1)\in T]|>\alpha.\]
Define tree $g$ with boolean output as $g(x):=1$ iff $g'(x)\in T$; note this just amounts to changing the labels of the leaves of $g'$. Now the left-hand side of the inequality above can be written as
\[|\mathbb P[g(R)=1]-\mathbb P[g(B_1)=1]|\]
and this contradicts the result for trees with boolean outputs and concludes the proof of the lemma.
